%----------------------------------------------------------------------------------------
% Anschreiben — Universität Greifswald, Institut für Biochemie
% NFDI4Cat: Wiss. Mitarbeiter/in Biotechnologie & Enzymkatalyse (Job-ID 16133)
% Ansprechpartner: Dr. Mark Dörr
%----------------------------------------------------------------------------------------

\documentclass[11pt,a4paper]{letter}

\usepackage[utf8]{inputenc}
\usepackage[T1]{fontenc}
\usepackage{lmodern}
\usepackage[ngerman]{babel}
\usepackage[a4paper,margin=2.4cm]{geometry}
\usepackage{parskip}

\signature{Kundai Farai Sachikonye}
\address{Kundai Farai Sachikonye \\ Zeltnerstraße 30 \\ 90443 Nürnberg \\ (+49) 1626386344 \\ kundai.sachikonye@bitspark.com}

\begin{document}

\begin{letter}{%
  Universität Greifswald \\
  Institut für Biochemie \\
  Herrn Dr. Mark Dörr \\
  Felix-Hausdorff-Straße 4 \\
  17489 Greifswald%
}

\opening{Bewerbung als wissenschaftlicher Mitarbeiter (w/m/d) im Projekt NFDI4Cat --- Job-ID 16133}

Sehr geehrter Herr Dr. Dörr,

mit großem Interesse bewerbe ich mich auf die Stelle im Teilprojekt der Universität
Greifswald innerhalb von NFDI4Cat. Die Ausschreibung beschreibt genau die
Schnittstelle, an der ich arbeite: \textbf{semantische Datenstandards und Ontologien
für maschinell gesteuerte, biokatalytische Prozesse} --- verbunden mit Laborautomation
und Closed-Loop-Experimentierung. Selten passt ein Profil so unmittelbar.

\textbf{Semantische Repräsentation und FAIR-Daten.} Ich entwerfe seit Jahren Ontologien,
Taxonomien und domänenspezifische Sprachen (DSLs) mitsamt ihren Compilern. Mehrere
meiner Projekte sind im Kern Aufgaben der semantischen Datenrepräsentation: die
strukturierte, eindeutige und maschinenlesbare Beschreibung wissenschaftlicher
Sachverhalte. Mit OWL/RDF, der Modellierung von Taxonomien und der Ausrichtung an den
FAIR-Prinzipien bin ich vertraut und arbeite täglich mit Python, JavaScript, Git,
Docker und agentischer Programmierung.

\textbf{Domänenwissen Enzymkatalyse.} Anders als rein informatische Bewerber bringe ich
ein fundiertes Verständnis der Chemie mit, die diese Daten beschreiben sollen. Ich habe
einen umfangreichen Forschungskorpus zur Enzymkatalyse verfasst --- u.\,a. zum
Mechanismus und zur Wechselzahl der Carboanhydrase sowie zu plattform-invarianten
Datenmodellen für die Massenspektrometrie. Datenstandards für biokatalytische Daten
sollten von jemandem mitgestaltet werden, der die Kinetik und die Mechanismen
tatsächlich versteht; diese Verbindung aus Domäne und Informatik biete ich.

\textbf{Agentik und Closed-Loop.} Mein Werkzeug \texttt{purpose} (Rust) ist eine
agenten-abfragbare Wissensschicht über wissenschaftlichen Projektkorpora --- genau das
Muster aus agentischer Programmierung und automatisierter, iterativer
Forschungsdatengenerierung, das die Stelle adressiert. Closed-Loop-Experimentierung,
also der Kreislauf aus Vorhersage, Messung und Modellanpassung, ist ein roter Faden
meiner Arbeit.

\textbf{Technisches Profil.} Fortgeschrittene Kenntnisse in Python und JavaScript, in
Container-Technologien (Docker, Kubernetes), in Datenbanken sowie in Git und
Entwicklerwerkzeugen; Erfahrung mit KI/Machine Learning (PyTorch) und mit logikbasiertem
Reasoning. Meine Englischkenntnisse sind auf muttersprachlichem Niveau (C2).

Einen formal qualifizierenden Hochschulabschluss (Master) im computergestützten Bereich
bringe ich mit; meine darüber hinausgehende, eigenständige Forschung dokumentiert das
geforderte überdurchschnittliche wissenschaftliche Niveau. Als Inhaber einer
unbefristeten Niederlassungserlaubnis in Deutschland verfüge ich über eine
uneingeschränkte Arbeitserlaubnis; eine Sponsoring-Regelung ist nicht erforderlich.

Über die Gelegenheit, das Vorhaben und das Greifswalder Teilprojekt in einem Gespräch
näher kennenzulernen, würde ich mich sehr freuen.

\closing{Mit freundlichen Grüßen}

\end{letter}
\end{document}
